\chapter{緒論}

\section{研究背景}

請在此說明研究主題的背景脈絡、應用場域與重要性。背景段落應協助讀者理解為何此問題值得研究，以及本論文與既有研究或實務需求之間的關係。

\section{研究動機}

請描述目前方法、系統、理論或實務流程中的不足之處。研究動機應明確連結到後續的研究問題與方法設計，避免只停留在一般性的敘述。

\section{研究目的與問題}

請列出本論文欲回答的核心問題。若研究包含多個子問題，可使用條列方式整理：

\begin{itemize}
  \item 研究問題一：請填入問題描述。
  \item 研究問題二：請填入問題描述。
  \item 研究問題三：請填入問題描述。
\end{itemize}

\section{研究範圍與限制}

請界定研究對象、資料來源、實驗條件、分析範圍與不處理的問題。清楚界定範圍可降低讀者對研究貢獻的誤解。

\section{論文架構}

本論文其餘章節安排如下：第二章回顧相關文獻；第三章說明研究方法；第四章呈現結果與討論；第五章提出結論、貢獻與未來研究方向。
